% from ega3-6-fr.tex
\oldpage[III]{174}
(qui est un complexe de $\mathscr O_{Y'}$-modules localement libres de type
fini), un isomorphisme
\[
  \mathscr H^\bullet
  (f',\mathscr P'_\bullet\otimes_{Y'}\mathscr Q'_\bullet)
  \xrightarrow{\sim}
  \mathscr H_\bullet
  (\mathscr L'_\bullet\otimes_{Y'}\mathscr Q'_\bullet)
  \tag{6.10.5.2}\label{III.6.10.5.2.fr}
\]
de $\partial$-foncteurs en le complexe $\mathscr Q'_\bullet$ de
$\mathscr O_{Y'}$-Modules quasi-cohérents, limité inférieurement, de sorte que
le diagramme
\[
\xymatrix@C=25pt@R=30pt{
  \mathscr H^\bullet
  (f,\mathscr P_\bullet\otimes_Y\mathscr Q_\bullet)
  \ar[r]^-{\sim}\ar[d]
  &
  \mathscr H_\bullet
  (\mathscr L_\bullet\otimes_Y\mathscr Q_\bullet)
  \ar[d]
  \\
  \mathscr H^\bullet
  (f',\mathscr P'_\bullet\otimes_{Y'}u^*(\mathscr Q_\bullet))
  \ar[r]^-{\sim}
  &
  \mathscr H_\bullet
  (\mathscr L'_\bullet\otimes_{Y'}u^*(\mathscr Q_\bullet))
}
\tag{6.10.5.3}\label{III.6.10.5.3.fr}
\]
soit commutatif.
\end{theorem}

\begin{proof}
L'application de (6.10.1) donne d'abord un complexe limité inférieurement
(6.10.2) $\mathscr R_\bullet$ de $\mathscr O_Y$-Modules quasi-cohérents et
$Y$-libres (6.10.3, (ii)) et un isomorphisme
\[
  \mathscr H^\bullet
  (f,\mathscr P_\bullet\otimes_Y\mathscr Q_\bullet)
  \xrightarrow{\sim}
  \mathscr H_\bullet
  (\mathscr R_\bullet\otimes_Y\mathscr Q_\bullet)
  \tag{6.10.5.4}\label{III.6.10.5.4.fr}
\]
de $\partial$-foncteurs en $\mathscr Q_\bullet$, mais \emph{a priori} les
termes de $\mathscr R_\bullet$ ne sont pas nécessairement des
$\mathscr O_Y$-Modules \emph{de type fini}. Mais si l'on applique (6.10.5.4) au cas
où $\mathscr Q_\bullet$ est un complexe réduit à un seul terme
$\mathscr O_Y$, on voit que $\mathscr H_\bullet(\mathscr R_\bullet)$ est
isomorphe à $\mathscr H^\bullet(f,\mathscr P_\bullet)$, et est par suite formé
de $\mathscr O_Y$-Modules \emph{cohérents} (6.2.5). On sait alors (\textbf{0}, 11.9.2) qu'il
existe un complexe $\mathscr L_\bullet$ limité inférieurement, formé de
$\mathscr O_Y$-Modules associés à des $A$-modules libres \emph{de type fini}, et un
homomorphisme $\mathscr L_\bullet\longrightarrow\mathscr R_\bullet$,
tels que l'homomorphisme correspondant pour l'homologie,
$\mathscr H_\bullet(\mathscr L_\bullet)\longrightarrow\mathscr H_\bullet(\mathscr R_\bullet)$
soit bijectif; d'où l'isomorphisme (6.10.5.1), en vertu de (6.10.3, (ii)).
Les autres assertions de (6.10.5) résultent de (6.10.1) et (6.10.3, (ii))
lorsque $Y'$ est \emph{affine}; dans le cas général, il suffit de vérifier que
lorsqu'on considère un recouvrement $(V_\alpha)$ de $Y'$ par des ouverts
affines, et l'isomorphisme correspondant (6.10.5.2) relatif à chacun des
$V_\alpha$, les restrictions à un ouvert affine
$W\subset V_\alpha\cap V_\beta$ des isomorphismes correspondant à
$V_\alpha$ et à $V_\beta$ coïncident avec l'isomorphisme correspondant à
$W$, ce qui résulte de la commutativité du diagramme (6.10.1.2) appliqué aux
injections canoniques $W\to V_\alpha$ et $W\to V_\beta$.
\end{proof}

\begin{remark}[6.10.6]
Dans les chapitres suivants, nous appliquerons surtout (6.10.5) au cas où
$\mathscr P_\bullet$ est réduit à un seul $\mathscr O_X$-Module cohérent
$\mathscr F$, plat sur $Y$. Comme on a alors
$\mathscr H_n(\mathscr L_\bullet)=\mathscr H^{-n}(f,\mathscr F)=\mathrm R^{-n}f_*(\mathscr F)$ (6.2.1),
on voit que les $\mathscr H_n(\mathscr L_\bullet)$ sont nuls pour $n>0$;
nous verrons plus loin (7.7.12, (i)) qu'on peut alors supposer que
$\mathscr L_\bullet$ \emph{n'a que des termes de degrés} $\leqslant0$ (donc en
nombre \emph{fini}), à condition de remplacer l'hypothèse que les $\mathscr L_i$
sont associés à des $A$-modules libres de type fini par celle que les
$\mathscr L_i$ sont \emph{localement libres} de type fini.

Le complexe $\mathscr L_\bullet$ correspondant à un tel
$\mathscr O_X$-Module $\mathscr F$ ne paraît posséder aucune
